\origpage{48}{N. Robertson, P.D. Seymour / Journal of Combinatorial Theory, Series B 89 (2003) 43--76}
\textbf{2.1.} \textit{For any $\theta \ge 1$, let $\mathscr{D}$ be a class of designs which is $\theta$-pervasive in a graph $G$. Then $G$ has a tree-decomposition $(T, \tau)$ such that for each $t \in V(T)$, the design of $t$ in $(T, \tau)$ belongs to $\mathscr{D}^{3\theta-2} \cup \mathscr{R}_{4\theta-3}$.}

Let $(H, M)$ be a design. A graph $H'$ is a \textit{torso} of $(H, M)$ if $V(H) = V(H')$, $H$ is a subgraph of $H'$, and for every $e \in E(H') \setminus E(H)$ there exists $X \in M$ including the ends of $e$. It is easy to see the following.

\textbf{2.2.} \textit{Let $\mathscr{D}$ be a class of designs, and let $\mathscr{D}'$ be the class of torsos of members of $\mathscr{D}$. If $G$ is a graph with a tree-decomposition $(T, \tau)$ such that $\mathscr{D}$ contains the design of $t$ in $(T, \tau)$ for each $t \in V(T)$, then $G$ is a tree-structure over $\mathscr{D}'$.}

\fix{\textbf{Proof.} For each edge $tt'$ of $T$, let $K_{tt'}$ be a complete graph with $V(K_{tt'}) = V(\tau(t) \cap \tau(t'))$, chosen so that no two of these complete graphs have a common edge, and none has an edge in common with $G$. For each subtree $S$ of $T$, let $G_S$ be the graph with vertex set $\cup(V(\tau(t)) \colon t \in V(S))$ and with edge set
\begin{equation*}
\cup(E(\tau(t)) \colon t \in V(S)) \cup \cup(E(K_{tt'}) \colon tt' \in E(T), t \in V(S), t' \notin V(S)).
\end{equation*}
Then $G_T = G$, since $\cup(\tau(t) \colon t \in V(T)) = G$. Hence it suffices to show that $G_S$ is a tree-structure over $\mathscr{D}'$ for every subtree $S$ of $T$, and we prove this by induction on $|V(S)|$.

Suppose first that $V(S) = \{t_0\}$, and let $t_1, \ldots, t_k$ be the neighbours of $t_0$ in $T$. Then $V(G_S) = V(\tau(t_0))$, $\tau(t_0)$ is a subgraph of $G_S$, and every other edge of $G_S$ is an edge of some $K_{t_0 t_i}$, and so its ends belong to $V(\tau(t_0) \cap \tau(t_i))$. Thus $G_S$ is a torso of the design of $t_0$ in $(T, \tau)$, and so $G_S \in \mathscr{D}'$.

Now suppose that $|V(S)| \ge 2$. Let $t_0$ be a vertex of $S$ with degree one in $S$, let $t_1$ be its neighbour in $S$, and let $S' = S \setminus t_0$. Let $H_1$ be the graph obtained from $G_{\{t_0\}}$ by deleting the edges of $K_{t_0 t_1}$, and let $H_2$ be obtained from $G_{S'}$ in the same way. We claim that $(H_1, H_2)$ is a separation of $G_S$ with $V(H_1 \cap H_2) = V(K_{t_0 t_1})$. Every neighbour $t'$ of $t_0$ in $T$ with $t' \ne t_1$ is not in $V(S)$, since $S$ is a subtree and $t_0$ has degree one in $S$. Hence the edges $tt'$ of $T$ with $t \in V(S)$ and $t' \notin V(S)$ are the edges $t_0 t'$ with $t' \ne t_1$, and the edges $tt'$ with $t \in V(S')$, $t' \notin V(S')$ and $t' \ne t_0$. It follows that $H_1 \cup H_2 = G_S$. No edge of $\tau(t_0)$ belongs to $\tau(t)$ with $t \ne t_0$, and the complete graphs $K_{t_0 t'}$ $(t' \ne t_1)$ appear only in $H_1$, so $H_1$ and $H_2$ have no common edges. Finally, let $t \in V(S')$. The path of $T$ between $t_0$ and $t$ lies in $S$, and so it passes through $t_1$. Hence $\tau(t_0) \cap \tau(t) \subseteq \tau(t_1)$, and so $V(\tau(t_0)) \cap V(\tau(t)) \subseteq V(\tau(t_0) \cap \tau(t_1))$. Since $V(\tau(t_1)) \subseteq V(H_2)$, it follows that $V(H_1 \cap H_2) = V(\tau(t_0) \cap \tau(t_1)) = V(K_{t_0 t_1})$. This proves our claim. Since $K_{t_0 t_1}$ has no edges in common with $H_1$ and $H_2$, and $H_1 \cup K_{t_0 t_1} = G_{\{t_0\}}$ and $H_2 \cup K_{t_0 t_1} = G_{S'}$, it follows that $G_S$ is the clique-sum of $G_{\{t_0\}}$ and $G_{S'}$. From the inductive hypothesis, $G_{\{t_0\}}$ and $G_{S'}$ are tree-structures over $\mathscr{D}'$, and hence so is $G_S$. This proves 2.2. $\square$}\fixnote{fix:2.2-proof-added}{The paper says only that 2.2 is easy to see. The argument above is the one used in the formalization.}

Now for $n \ge 0$, if $\mathscr{D}$ is a class of designs then every torso of a member of $\mathscr{D}^n$ is a $(\le n)$-vertex extension of a torso of a member of $\mathscr{D}$\fix{, provided that $Z \cap V(H) = \emptyset$ in the definition of $\mathscr{D}^n$. In general, such a torso $G$ contains a torso $G'$ of a member of $\mathscr{D}$, and $G \setminus X = G' \setminus X$ for some $X \subseteq V(G)$ with $|X| \le n$}.\fixnote{fix:2-torso}{As printed, the sentence fails when $Z$ meets $V(H)$, which page 47 allows, because then a torso of $(H', M')$ can have edges with both ends in $Z \cap V(H)$. For example, let $H$ have two vertices $a, b$ and no edges, and let $M = \emptyset$, $H' = H$, $M' = \{\{a, b\}\}$, $Z = \{a, b\}$ and $n = 2$. The graph with the one edge $ab$ is a torso of $(H', M')$. The only torso of $(H, M)$ is $H$, and $H$ is not obtained from that graph by deleting vertices. With the weaker form inserted here, 2.3 holds with the same weakening in its second case, and the proof of 1.3 still goes through, because deleting at most $n$ vertices from an outgrowth by $\le d$ $r$-rings gives an outgrowth by $\le d$ $(n+1)r$-rings. So 1.3 holds as printed, with a larger $r$ than the proof's $2\rho + 1$.} From 2.1 and 2.2 we deduce the following.

\textbf{2.3.} \textit{For any $\theta \ge 1$, let $\mathscr{D}$ be a class of designs which is $\theta$-pervasive in a graph $G$. Then $G$ may be constructed by clique-sums starting from graphs $G'$ such that either}

\hspace*{12pt}$\bullet$ \textit{$|V(G')| \le 4\theta - 3$, or}

\hspace*{12pt}$\bullet$ \textit{$G'$ \del{is a $(\le 3\theta - 2)$-vertex extension of} \fix{contains} a torso \fix{$G''$} of a member of $\mathscr{D}$\fix{, and $G' \setminus X = G'' \setminus X$ for some $X \subseteq V(G')$ with $|X| \le 3\theta - 2$}.}\fixnote{fix:2.3}{A gap in the proof. The proof of 2.3 uses the sentence before it, which holds only in the weaker form given there, so it proves 2.3 only with the weaker second case given here. We do not know whether 2.3 holds as printed. Its only use is in the proof of 1.3, on page 52, and the weaker form is enough there.}

We shall use 2.3 to derive 1.3 from a theorem that a certain class of designs is $\theta$-pervasive.

\sectionhead{3. Surfaces, societies and segregations}

In this paper, by a \textit{surface} we mean a non-null compact connected 2-manifold without boundary. An \textit{$O$-arc} in a surface $\Sigma$ is a subset $F \subseteq \Sigma$ homeomorphic to a circle. \textit{Open} and \textit{closed discs} in $\Sigma$ are defined in the natural way. For $X \subseteq \Sigma$, its \textit{closure} is denoted by $\bar{X}$, and $\bar{X} \cap \overline{\Sigma \setminus X}$ is denoted by $bd(X)$. If $F \subseteq \Sigma$ is an $O$-arc and $X \subseteq F$ is finite then $F$ induces two cyclic permutations on $X$, called the \textit{natural orders of $X$ from $F$}.

A \textit{society} is a pair $(A, \Omega)$, where $A$ is a graph and $\Omega$ is a cyclic permutation of a subset (denoted by $\bar{\Omega}$) of $V(A)$. A \textit{segregation} of $G$ is a set $\mathscr{S}$ of societies such that

\hspace*{12pt}$\bullet$ $A \subseteq G$ for every $(A, \Omega) \in \mathscr{S}$, and $\cup(A \colon (A, \Omega) \in \mathscr{S}) = G$,

\hspace*{12pt}$\bullet$ $V(A \cap A') \subseteq \bar{\Omega} \cap \bar{\Omega}'$ and $E(A \cap A') = \emptyset$ for all distinct $(A, \Omega)$, $(A', \Omega') \in \mathscr{S}$.

We write $V(\mathscr{S}) = \cup(\bar{\Omega} \colon (A, \Omega) \in \mathscr{S})$.

Let $\Sigma$ be a surface, and $\mathscr{S} = \{(A_1, \Omega_1), \ldots, (A_k, \Omega_k)\}$ a segregation of $G$. It is convenient always to assume (as we may) that $S \cap V(\mathscr{S}) = \emptyset$. An \textit{arrangement}
